\chapter{Features}
\label{chap:features}

This template provides an out-of-box experience for turning a written article into a journal issue, which can be used for both print and online versions.

\section{Compatibility}

This template is compatible with the template in \cite{article-template}.
This means that you can write your article using \cite{article-template} and then directly copy and paste the content into this template.
The only modification you need to make is to change the ``section'' commands to ``chapter'' commands.

\section{Two formats}

This template provides two formats for the journal issue: a print format and an online format.

\begin{description}
  \item[Print]
    This format is designed for printing the journal issue.
    It uses a two-sided layout, which means that the headers are different on odd and even pages, the margins are adjusted accordingly, and chapters start on odd pages.
    The front and back covers are not included in the document, for the convenience of printing and binding.
  
  \item[Online]
    This format is designed for online reading of the journal issue.
    It uses a one-sided layout.
    The front and back covers are directly included in the document.
\end{description}

You can run \verb|./build.sh| to build both formats, and the output files will be saved in the \verb|target/| directory.
See \Cref{chap:usage} for more details on how to build the journal issue.

\section{Cover design}

The covers of the journal issue is designed with Inkscape, which is a free and open-source vector graphics editor.
The source files of the covers are included in the \verb|resources/| directory, and you can modify them directly (e.g., update the volume and issue numbers).

\verb|./resources/export.sh| is a script that can be used to export the covers from \verb|.svg| format to \verb|.png| format, which is the format used in the journal issue.
This is because GitHub actions do not support Inkscape, so we cannot directly use the \verb|.svg| files in the journal issue.

\section{GitHub actions}

This template uses GitHub actions to automate the building and publishing of the journal issue.
When you push changes to the \verb|main| branch, the GitHub action will automatically build both formats of the journal issue and publishes them onto GitHub Pages.
The published journal issue can be accessed at \url{\docpublisherwebsite/<repo-name>}.
